\section{ADR-012: Audio to MIDI per stem}
\textbf{Status:} accepted and implemented (Future work, FR-MIDI-01).

\textbf{Context.} Musicians who learn or arrange a cover want the notes of a
part, not only its sound. The separated stems are clean input for automatic
transcription, and the beat grid of the analysis (ADR-010) gives the bars.

\textbf{Pitched stems: Basic Pitch on ONNX Runtime.} Spotify's Basic Pitch
(ICASSP 2022) is a small polyphonic note model trained across instruments;
its code and weights are Apache-2.0. The \texttt{basic-pitch} package cannot
be installed: its last release (0.4.0, 2024) requires TensorFlow
\textless\,2.15.1 on Windows with Python 3.12, which has no wheels for it. So
the engine ports the inference and the note decoding (about 300 lines, with
Spotify's notice and a list of changes) to numpy, and runs the unmodified
ONNX model (230\,KB) with ONNX Runtime (MIT), which the separation runtime
already installed and which is now a base dependency (15\,MB wheel for
cp312). The model ships in the engine package with its \texttt{LICENSE} and
\texttt{NOTICE}: the one bundled model, kept because it is tiny and
permissively licensed; separation models stay a user download. Pitch bends
are not written. Each stem has a profile: a frequency range (vocals 65--1400\,Hz,
bass 30--420\,Hz, guitar 75--1500\,Hz), a General MIDI program, and for
vocals and bass a monophonic pass (the louder of two overlapping notes wins),
which also drops the octave overtones that the model reports on loud notes.

\textbf{Drums: signal processing only.} No permissively licensed drum
transcription model was found to bundle. The drum stem's power spectrum is
summed in three bands (kick 30--150\,Hz, snare 150--4000\,Hz, hi-hat
6--11\,kHz); a candidate hit is a sharp rise of any band, and a drum sounds
when its band gained at least a tenth of its typical hit and 40\,\% of the
strongest band's share at that moment, because a kick's beater click leaks
into the snare band. Notes go to channel 10 (36, 38, 42). Everything is
librosa (ISC) and numpy (BSD).

\textbf{Where it runs.} On CPU threads of the engine process (4 intra-op
threads), not in the GPU worker: the model is small, the worker stays free for
separation, and the work needs no PyTorch. On the target laptop
(i5-11400H) a 4-minute stem takes about 3.5\,s (2.2\,s of model and 1.2\,s of
note decoding), 66 times faster than real time; all 12 threads gain only
5\,\%.

\textbf{Tempo map and file.} The files are Standard MIDI Files, format 1,
960 ticks per quarter note: a conductor track with the song title, time
signatures and tempo changes, then the stem's track with its name and
program. The tempo map is fitted to the downbeats of the analysis with the
algorithm of the DAW session export (ADR-011): one tempo for steady songs,
one per bar otherwise, and a pickup bar before the first downbeat; without a
grid, one 120\,BPM tempo. Note times in seconds become ticks through the same
map, so a DAW that follows the file's tempo plays each note where it sounds in
the stem. Until both changes are merged the algorithm lives in
\texttt{midi/tempo.py}; afterwards one copy remains.

\textbf{Pipeline step and API.} \texttt{POST /items/\{id\}/midi} takes the
stem ids; the engine plans the work on the database, transcribes each stem on
a thread and emits \texttt{midi.progress} events (stem, index, total,
fraction). The step \texttt{midi} has its own cache key: the stem's asset
SHA-256, the transcription version and a digest of the tempo map. The MIDI
file is kept in the asset store and copied to the job's output folder as
\texttt{<title> - <Stem>.mid}, next to the stem exports, so asking again
never transcribes a stem twice.

\textbf{Consequences.} Transcription is automatic: polyphonic stems keep some
octave overtones, fast or quiet notes can be missed, and drums other than
kick, snare and hi-hat are left out or merged. The bundled model adds 230\,KB
and a licence notice; the README and the notices now say that separation
weights, not all weights, are never bundled.
